\documentclass{situationist-preprint}

\title{Ockhamism and the Fixity of the Past}
\msid{REL-05-b}
\msfamily{REL-05}
\msdiscipline{Philosophy of Religion}
\mstype{Companion manuscript}
\mskeywords{Ockhamism; soft facts; fixity of the past; branching time; divine foreknowledge; future contingents}

\begin{document}
\maketitle

\begin{abstract}
The Ockhamist reconciles foreknowledge with freedom by classifying God's past beliefs about future choices as soft facts, which are not fixed. REL-05-a sets this strategy aside and grants the fixity of the past without restriction. We argue that the concession costs nothing and that Ockhamism, in branching time with moment-relative settledness, is not an alternative to the account of REL-05-a but a confused statement of it. Every fact about a past moment is settled at every later moment; a divine belief at an earlier moment about a contingent choice would have to be settled before it was, which branching excludes. What the Ockhamist calls a soft fact is a tenseless truth relative to a history, and that is exactly the object of timeless knowledge.
\end{abstract}

\section{The Ockhamist strategy}

Pike's argument holds that if God believed at $t_0$ that Jones would mow his lawn at $t_1$, then Jones could refrain at $t_1$ only if he could make it the case that God held a false belief, or that God did not hold the belief God held, and neither is within anyone's power \citep{Pike1965}. The Ockhamist reply, revived by \citet{Adams1967} and developed by \citet{Plantinga1986}, distinguishes hard facts, which are genuinely about the past, from soft facts, which are about the past only in grammatical form because they depend on what happens later. \enquote{Jones mowed at $t_1$, 80 years after $t_0$} is about $t_0$ only in form. If God's belief at $t_0$ is soft, it is not fixed, and Jones has the power to act in such a way that God would not have held it. The reply has met two difficulties: no criterion of softness has commanded agreement, and divine belief seems to have a hard core, since God's believing something is a state of God at $t_0$ \citep{Zagzebski1991}.

REL-05-a reaches compatibility by a different route. It grants that a temporally located divine belief before the choice would be fixed (its Assumption~2.7), and concludes from the agent's ability to do otherwise that God's knowledge is not such a fact (its Corollary~3.4): it is timeless, and depends counterfactually on the choice. We show that the framework of tenses and times that the account presupposes leaves the Ockhamist no further room.

\section{Settledness and the past}

We work in the branching-time framework of \citet{Thomason1970} and \citet{Belnap2001}, combined with a tenseless view of time. Moments are partially ordered; a history is a maximal chain of moments; a sentence is evaluated at a moment $m$ relative to a history $h$ through $m$. A tenseless sentence such as \enquote{Jones mows at $t_1$} is true relative to a history just in case the history contains a moment at $t_1$ at which Jones mows.

\begin{definition}\label{d:settled}
A proposition is \emph{settled} at $m$ if it has the same truth value relative to every history through $m$. A future contingent at $m$ is a proposition about a later time that is not settled at $m$.
\end{definition}

Following TSM-12, a future contingent is neither true nor false at $m$; it becomes true or false at the moment at which it is settled. Following TSM-01 and TSM-03, this is not a claim that what is future lacks existence: there is no privileged present, and later moments are elsewhere in time. Openness is a modal fact about a moment, namely that more than one history passes through it.

\begin{proposition}\label{p:fixed}
Every proposition that is settled at $m$ is settled at every later moment $m'$. In particular, everything about moments up to $m$ that is settled at $m$ is fixed at $m'$.
\end{proposition}

\begin{proof}
Every history through $m'$ passes through $m$, since histories are maximal chains and $m < m'$. A proposition with the same value on all histories through $m$ therefore has that value on all histories through $m'$.
\end{proof}

Proposition~\ref{p:fixed} is the fixity of the past in this framework, and it holds without exception. The Ockhamist must therefore deny not the fixity of settled facts but that the divine belief is settled at $t_0$.

\begin{proposition}\label{p:nobelief}
Let $c$ be a choice at $t_1$ that is a future contingent at a moment $m_0$ at $t_0$. There is no settled fact at $m_0$ that God knows at $t_0$ that the agent takes $c$.
\end{proposition}

\begin{proof}
Suppose it were settled at $m_0$ that God knows at $t_0$ that the agent takes $c$. Since divine knowledge is factive (REL-05-a, Assumption~2.3), what God knows is true relative to every history through $m_0$, so the agent takes $c$ on every such history, and the choice is settled at $m_0$, contrary to the hypothesis.
\end{proof}

The Ockhamist can now say only that the divine belief at $t_0$ obtains relative to some histories through $m_0$ and not others: relative to those in which the agent takes $c$. But a fact about $t_0$ that holds relative to a history only in virtue of what that history contains at $t_1$ is not a fact about $t_0$ at all, in the sense in which a state of God at $t_0$ would be. It is a tenseless truth relative to a history, of the form \enquote{the agent takes $c$ at $t_1$}, given a temporal tag. This is the soft fact, correctly described, and it is precisely what REL-05-a's timeless knowledge has as its object: the truth of a proposition relative to the actual history, not its truth at a moment before the choice.

\section{What remains of Ockhamism}

Two of the Ockhamist's claims survive, and both belong to the account of REL-05-a. The first is that the agent's power over \enquote{God's belief} is counterfactual: had she chosen otherwise, God would have known otherwise. This is the first half of the family's conclusion, Corollary~3.3 of REL-05-a, and it is Plantinga's counterfactual power over the past with the word \enquote{past} removed. The second is that this power does not extend to anything that is genuinely about an earlier time. By Proposition~\ref{p:fixed}, nothing settled at an earlier moment is within the agent's power; by Proposition~\ref{p:nobelief}, God's knowledge of her contingent choice was never settled at an earlier moment, so it was never among the things fixed.

What does not survive is the claim that God has, at $t_0$, beliefs about future contingents. That claim would make the divine belief a hard fact whose content is unsettled, and Proposition~\ref{p:nobelief} shows that there is no such thing.

\begin{objection}
The Ockhamist may adopt a thin red line: a privileged actual history, relative to which future contingents are true at earlier moments. God's belief at $t_0$ is then true relative to that line, and soft.
\end{objection}

\begin{reply}
On a thin red line, the actual history's later contents fix what is true at $t_0$, and the belief is either settled at $m_0$, in which case Proposition~\ref{p:nobelief} applies, or holds only relative to the line, in which case it is the tenseless truth just described. \citet{BelnapGreen1994} argue that the thin red line cannot be integrated with branching without making the alternatives idle; on our framework it adds nothing that the tenseless truth relative to the actual history does not already supply.
\end{reply}

\section{Conclusion}

In branching time with moment-relative settledness, the past is fixed without exception, and God's knowledge of a free choice was never part of it. The Ockhamist's soft facts are tenseless truths relative to a history; known timelessly, they depend counterfactually on the choice, and the choice does not depend on them. Ockhamism is the account of REL-05-a misdescribed as a thesis about the past.

\begin{thebibliography}{99}\small
\bibitem[Adams(1967)]{Adams1967} Adams, M.~M. (1967). Is the existence of God a \enquote{hard} fact? \emph{The Philosophical Review}, 76(4), 492--503.
\bibitem[Belnap et~al.(2001)]{Belnap2001} Belnap, N., Perloff, M. and Xu, M. (2001). \emph{Facing the Future: Agents and Choices in Our Indeterminist World}. Oxford University Press.
\bibitem[Belnap and Green(1994)]{BelnapGreen1994} Belnap, N. and Green, M. (1994). Indeterminism and the thin red line. \emph{Philosophical Perspectives}, 8, 365--388.
\bibitem[Pike(1965)]{Pike1965} Pike, N. (1965). Divine omniscience and voluntary action. \emph{The Philosophical Review}, 74(1), 27--46.
\bibitem[Plantinga(1986)]{Plantinga1986} Plantinga, A. (1986). On Ockham's way out. \emph{Faith and Philosophy}, 3(3), 235--269.
\bibitem[Thomason(1970)]{Thomason1970} Thomason, R.~H. (1970). Indeterminist time and truth-value gaps. \emph{Theoria}, 36(3), 264--281.
\bibitem[Zagzebski(1991)]{Zagzebski1991} Zagzebski, L.~T. (1991). \emph{The Dilemma of Freedom and Foreknowledge}. Oxford University Press.
\end{thebibliography}

\end{document}
